% from ega4-8-fr.tex
\oldpage[IV-3]{24}
  exacte, il faut et il suffit qu'il existe $\mu\geq\lambda$ tel que la suite
  $\mathcal F_\mu\xrightarrow{f_\mu}\mathcal G_\mu
  \xrightarrow{g_\mu}\mathcal H_\mu$ soit exacte.

  \item Soit $f_\lambda:\mathcal F_\lambda\to\mathcal G_\lambda$ un
  homomorphisme de $\mathcal O_{S_\lambda}$-Modules quasi-cohérents tel que
  $\operatorname{Ker}f_\lambda$ soit de type fini. Pour que
  $f=u_\lambda^*(f_\lambda):\mathcal F\to\mathcal G$ soit injectif, il faut et
  il suffit qu'il existe $\mu\geq\lambda$ tel que
  $f_\mu=u_{\lambda\mu}^*(f_\lambda):\mathcal F_\mu\to\mathcal G_\mu$ le soit.
\end{enumerate}
\end{corollary}

\begin{proof}
\begin{enumerate}[label=(\roman*)]
  \item Compte tenu de (8.3.8, (ii)), notons que, par platitude,
  $\operatorname{Im}f$ et $\operatorname{Ker}g$ (resp.
  $\operatorname{Im}f_\mu$ et $\operatorname{Ker}g_\mu$ pour
  $\mu\geq\lambda$) sont les images réciproques de
  $\operatorname{Im}f_\lambda$ et $\operatorname{Ker}g_\lambda$
  $(0_{\mathrm I},6.7.2)$. Supposons que la suite
  $\mathcal F\xrightarrow{f}\mathcal G\xrightarrow{g}\mathcal H$ soit exacte.
  Puisque $\operatorname{Im}f_\lambda$ est de type fini, il existe
  $\mu\geq\lambda$ tel que le composé
  $\mathcal F_\mu\xrightarrow{f_\mu}\mathcal G_\mu
  \xrightarrow{g_\mu}\mathcal H_\mu$ soit nul, en vertu de (8.5.2, (i)).
  Changeant les notations, on peut donc déjà supposer que
  $g_\lambda\circ f_\lambda=0$. Alors comme l'homomorphisme
  $\mathcal F\to\operatorname{Ker}g$ est surjectif et que
  $\operatorname{Ker}g_\lambda$ est de type fini, il résulte de (8.5.7) qu'il
  existe $\mu\geq\lambda$ tel que l'homomorphisme
  $\mathcal F_\mu\to\operatorname{Ker}g_\mu$ soit surjectif, ce qui achève de
  prouver (i).

  \item L'assertion est le cas particulier de (i) appliqué à la suite
  $0\to\mathcal F_\lambda\to\mathcal G_\lambda$.
\end{enumerate}
\end{proof}

\begin{lemma}[8.5.9]
\label{IV.8.5.9-fr}
Supposons $S_0$ quasi-compact, $\mathcal F_\lambda$ quasi-cohérent de type
fini~; soient $\mathcal G'_\lambda$ et $\mathcal G''_\lambda$ deux quotients
quasi-cohérents de $\mathcal F_\lambda$, $\mathcal G'_\lambda$ étant en outre
supposé de présentation finie. Pour que $\mathcal G''$ soit quotient de
$\mathcal G'$, il faut et il suffit qu'il existe $\mu\geq\lambda$ tel que
$\mathcal G''_\mu$ soit quotient de $\mathcal G'_\mu$.
\end{lemma}

\begin{proof}
Par hypothèse, il y a deux homomorphismes surjectifs
$p'_\lambda:\mathcal F_\lambda\to\mathcal G'_\lambda$,
$p''_\lambda:\mathcal F_\lambda\to\mathcal G''_\lambda$~; en vertu de
l'exactitude à droite de $u_\lambda^*$ et de $u_{\lambda\mu}^*$,
$p'=u_\lambda^*(p'_\lambda)$, $p''=u_\lambda^*(p''_\lambda)$,
$p'_\mu=u_{\lambda\mu}^*(p'_\lambda)$,
$p''_\mu=u_{\lambda\mu}^*(p''_\lambda)$ sont aussi surjectifs~; en outre, s'il
existe un homomorphisme $f:\mathcal G'\to\mathcal G''$ (resp.
$f_\mu:\mathcal G'_\mu\to\mathcal G''_\mu$) tel que $p''=f\circ p'$ (resp.
$p''_\mu=f_\mu\circ p'_\mu$), cet homomorphisme est nécessairement unique, ce
qui montre que la question est locale sur $S_\lambda$, et qu'on peut par suite
($S_0$ étant quasi-compact et $L$ filtrant) supposer $S_\lambda$ affine, donc
quasi-séparé. Il est clair que la condition de l'énoncé est suffisante.
Inversement, comme $\mathcal G'_\lambda$ est de présentation finie,
$S_\lambda$ quasi-compact et quasi-séparé, il résulte de (8.5.2, (i)) que s'il
existe un homomorphisme $f:\mathcal G'\to\mathcal G''$ tel que
$p''=f\circ p'$, il existe un $\mu\geq\lambda$ et un homomorphisme
$f_\mu:\mathcal G'_\mu\to\mathcal G''_\mu$ tels que
$f=u_\mu^*(f_\mu)$ et $p''_\mu=f_\mu\circ p'_\mu$, d'où le lemme.
\end{proof}

\begin{env}[8.5.10]
\label{IV.8.5.10-fr}
Dans la fin de ce numéro, pour tout Module quasi-cohérent $\mathcal F$ sur un
préschéma, notons $\mathfrak Q(\mathcal F)$ l'ensemble des Modules quotients de
$\mathcal F$ qui sont de présentation finie. Si $\mathcal F_\lambda$ est
quasi-cohérent et $\mathcal G_\lambda\in\mathfrak Q(\mathcal F_\lambda)$, il
résulte du fait que $u_{\lambda\mu}^*$ et $u_\lambda^*$ sont exacts à droite,
que l'on a $\mathcal G_\mu\in\mathfrak Q(\mathcal F_\mu)$ pour
$\mu\geq\lambda$ et $\mathcal G\in\mathfrak Q(\mathcal F)$~; il est clair que
$(\mathfrak Q(\mathcal F_\lambda),u_{\lambda\mu}^*)$ est un système inductif
d'ensembles, et que les
$u_\lambda^*:\mathfrak Q(\mathcal F_\lambda)\to\mathfrak Q(\mathcal F)$
forment un système inductif d'applications, d'où, en passant à la limite
inductive, une application canonique
\begin{equation}
  u_{\mathfrak Q}:\varinjlim\mathfrak Q(\mathcal F_\lambda)
  \longrightarrow\mathfrak Q(\mathcal F).
  \tag{8.5.10.1}\label{IV.8.5.10.1-fr}
\end{equation}

En outre, si $(\mathcal F'_\lambda)$ est une seconde famille de
$\mathcal O_{S_\lambda}$-Modules quasi-cohérents et si, pour tout $\lambda$,
$\mathcal F'_\lambda$ est un quotient de $\mathcal F_\lambda$, alors
$\mathcal F'$ est un quotient de $\mathcal F$ et l'on a un diagramme
commutatif

